% Brief (grammar: layered assurance stack with two rails; new addition, not
%   in the original figure list -- addendum: ch2's no-visual run at
%   Implementation Verification / the Arbiter. Source: foundling
%   whitepaper-foundlings/.../fig-anchor-verification-stack.tex, redrawn to
%   v2 roles; facts checked against Table tab:anchor-verification-stack and
%   its lead paragraph before adoption).
% Reader question: what does each verification layer actually show, and
%   where does a human, not a tool, have to assert the bridge holds?
% Claim: obligations descend three layers (protocol model, Rust source,
%   running binary); evidence returns upward; two irreducible gaps (the
%   specification gap, the deployment gap) are exactly where a human, not a
%   tool, asserts the bridge.
% Evidence: Table tab:anchor-verification-stack (ProVerif / Kani+tests /
%   conformance suite; the specification and deployment gaps named in its own
%   caption) [internal].
% Must distinguish: what a tool showed (the evidence rail) from what remains
%   asserted, never shown (the two gap markers) -- conflating them is the
%   false inference this figure exists to block.
% Grammar chosen: two-rail layered stack (nearest template: tpl-block-stack).
%   Rejected: repeating Table tab:anchor-verification-stack as a second table
%   -- the reader already has the table; the figure's job is the one fact the
%   table does not draw, that the gaps sit between specific layers on a
%   directed bridge, not scattered in a column.
% Five-second test: "three layers, obligations down the left, evidence up the
%   right, two amber gaps where only a human bridges the claim."
\begin{figure}[H]
\centering
\pdfigurehue{pdcobalt}%
\begin{tikzpicture}[pd figure,x=.9cm,y=1cm]
  \foreach \y/\title/\tool in {
    2.6/{protocol model}/{ProVerif},
    0/{Rust source}/{Kani + tests},
    -2.6/{running binary}/{conformance suite}} {
    \draw[pd hairline] (-1.3,\y) -- (3.9,\y);
    \node[pd title,anchor=west] at (-5.2,\y+.05) {\title};
    \node[pd label,anchor=west,text=pdinkmuted] at (-5.2,\y-.35) {\tool};
    \node[pd focus datum] at (3.9,\y) {};
  }

  \draw[pd arrow,line width=1.0pt] (-6.2,2.95) -- (-6.2,-2.95);
  \node[pd label,anchor=south,rotate=90,text width=1.7cm,align=center] at (-6.4,0) {obligations};
  \draw[pd focus arrow,line width=1.0pt] (4.65,-2.95) -- (4.65,2.95);
  \node[pd focus label,anchor=north,rotate=90,text width=1.7cm,align=center] at (4.85,0) {evidence};

  \foreach \y/\gap in {1.3/{specification gap},-1.3/{deployment gap}} {
    \draw[pd warn rule] (-3.7,\y) -- (3.9,\y);
    \node[pd warn datum,minimum size=4.8pt,inner sep=0pt] at (0,\y) {};
    \node[pd label,anchor=south,text width=3.2cm,align=center] at (0,\y+.22) {\gap};
  }
\end{tikzpicture}
\caption{The assurance case as a two-rail bridge, restating
Table~\ref{tab:anchor-verification-stack} as one directed shape.
Obligations descend from the protocol model to the Rust source to the
deployed binary; evidence from each tool's own run climbs back up the
right rail. The two amber gaps are the only places in the whole case where
a human, not a tool, asserts a layer's claim carries to the next --- that
the modeled protocol is this source, and that the inspected source is this
running binary. [internal]}
\label{fig:anchor-verification-stack}
\end{figure}
